\documentclass{beamer}
\usetheme{Madrid}
\usecolortheme{whale}
\title[Sustentación Proyecto]{Predicción de Deserción Estudiantil con Minería de Datos}
\subtitle{Máquina Virtual de Minería de Datos}
\author{Equipo de Investigación}
\date{\today}
\begin{document}
\begin{frame}
\titlepage
\end{frame}
\begin{frame}{1. Problema y Contexto}
\begin{itemize}
\item La deserción temprana en primer año genera impactos académicos y de gestión universitaria.
\item Objetivo: Evaluar si los modelos de minería de datos detectan oportunamente casos de riesgo.
\item Base de datos: Muestra didáctica de 500 registros con estructura inspirada en UCI \#697; no verificada como submuestra.
\item Distribución muestral: 5 casos positivos de abandono en 500 registros (1.0\% de prevalencia).
\end{itemize}
\end{frame}
\begin{frame}{2. Benchmark de Modelos en R}
Resultados en el conjunto de prueba independiente (150 registros, 3 positivos):
\begin{itemize}
\item \textbf{Regresión Logística:} Exactitud 98.00\% | Sensibilidad: 0\%.
\item \textbf{Árbol CART:} Exactitud 98.00\% | Sensibilidad: 0\%.
\item \textbf{Red Neuronal MLP:} Exactitud 98.00\% | Sensibilidad: 0\%.
\item \textbf{Baseline de Clase Mayoritaria:} 98.00\%.
\end{itemize}
\end{frame}
\begin{frame}{3. Diagnóstico Crítico y Recomendaciones}
\begin{itemize}
\item Los tres modelos igualan el baseline trivial de la clase mayoritaria (98.00\%).
\item Diagnóstico: La regla aprendida clasifica todas las instancias como no abandono ante el desbalance extremo de la muestra didáctica.
\item Recomendaciones: Aplicar técnicas de balanceo (SMOTE, submuestreo) y evaluar sobre el dataset completo de la fuente oficial.
\end{itemize}
\end{frame}
\end{document}
